\title{The Era of 1-bit LLMs: All Large Language Models are in 1.58 Bits}

\begin{abstract}
Advances like BitNet~\cite{bitnet} are ushering in a transformative phase for 1-bit Large Language Models (LLMs). In this study, we present a ternary-weight LLM called \textbf{\text{BitNet b1.58}}, where each model parameter is constrained to values in \{-1, 0, 1\}. This design achieves parity with full-precision Transformers (using FP16 or BF16), matching them in perplexity and downstream task effectiveness under equivalent model sizes and amounts of training data. At the same time, it dramatically improves efficiency regarding inference speed, memory usage, throughput, and power demands. More importantly, the introduction of the 1.58-bit LLM proposes a novel scaling relationship and training guideline for next-generation LLMs, emphasizing both cost efficiency and strong performance. Additionally, this direction lays the groundwork for alternative computational frameworks and paves the way for hardware architectures specifically tailored for 1-bit LLM workloads.
\end{abstract}

\section{The Era of 1-bit LLMs}

Recent developments in artificial intelligence have resulted in substantial expansion in both the scale and the capabilities of Large Language Models (LLMs). While these advances have led to impressive results across various natural language processing applications, the rapid growth in model parameters has introduced significant obstacles for real-world deployment, as well as heightened environmental and financial costs driven by substantial power requirements.

A promising strategy to mitigate these difficulties involves leveraging post-training quantization methods to generate low-bitwidth models suitable for inference~\cite{smoothquant, gptq, quip, quip_sharp}. By lowering the numerical precision of both parameters and intermediate activations, this approach can sharply decrease both memory footprints and compute overhead for LLMs. There has been an industry shift from standard 16-bit implementations toward ultra-low-bit variants, such as 4-bit models~\cite{gptq, awq}. Nonetheless, despite their practical adoption, post-training quantization methods remain fundamentally limited and do not achieve optimal performance.

Recent advancements in 1-bit model designs, exemplified by BitNet~\cite{bitnet}, have opened up new possibilities for reducing the operational costs associated with LLMs while preserving their efficacy. Standard LLMs rely on 16-bit floating-point representations (such as FP16 or BF16), with matrix multiplication accounting for most computational overhead. Consequently, the dominant source of computation cost arises from performing floating-point addition and multiplication. By contrast, BitNet performs matrix multiplication using only integer addition, drastically reducing the energy requirements associated with running LLMs. As power constraints often define the maximum computational throughput in modern hardware, these energy efficiencies can directly lead to accelerated processing.

Beyond computational expense, the overhead involved in transferring parameters from off-chip DRAM to on-chip accelerator memory (e.g., SRAM) during inference is also significant. Expanding SRAM to improve throughput has been proposed, but this approach incurs considerably higher costs relative to DRAM. Relative to their full-precision counterparts, 1-bit LLMs demand far less memory in terms of both storage capacity and bandwidth, which can notably decrease the expense and duration of loading model weights, thereby enabling swifter and more economical inference.

In this paper, we present a novel 1-bit LLM extension, referred to as \textbf{\text{BitNet b1.58}}, where all parameters are constrained to ternary values: \{-1, 0, 1\}. The addition of the 0 value differentiates this model from the original 1-bit BitNet, yielding a bit-width of 1.58 in binary encoding. \text{BitNet b1.58} retains all the core advantages of BitNet, including its innovative computation framework—essentially eliminating multiplication in matrix operations, facilitating aggressive hardware-level optimization. The model maintains the same energy profile as the 1-bit version and demonstrates marked improvements over FP16 LLMs in memory consumption, throughput, and latency. Moreover, \text{BitNet b1.58} introduces two further benefits. Firstly, it enhances expressive power by explicitly enabling feature filtering via the 0-valued weights, which bolsters the effectiveness of 1-bit LLMs. Secondly, experimental results demonstrate that, starting at a scale of 3B parameters and holding training configuration constant (model size, token count, etc.), \text{BitNet b1.58} can reach parity with FP16 models in both perplexity and downstream task metrics.

\section{BitNet b1.58}

\text{BitNet b1.58} builds upon the BitNet framework, a Transformer variant where the traditional \emph{nn.Linear} layers are substituted with \emph{BitLinear} layers. This model is initialized and trained from the ground up, employing weights quantized to 1.58 bits and activations quantized to 8 bits. This version features several modifications to the initial BitNet, which are outlined as follows.

\paragraph{Quantization Function.} To ensure that the weight values are restricted to -1, 0, or +1, we implement an \emph{absmean} quantization approach. Here, each entry in the weight matrix is divided by its mean absolute value, and subsequently rounded to the closest value among \{-1, 0, +1\}:

\begin{equation}
    \widetilde{W} = \mathrm{RoundClip}(\frac{W}{\gamma + \epsilon}, -1, 1),
\end{equation}

\begin{equation}
    \mathrm{RoundClip}(x, a, b) = \max(a, \min(b, \mathrm{round}(x))),
\end{equation}

\begin{equation}
    \gamma = \frac{1}{nm}\sum_{ij} |W_{ij}|.
\end{equation}

The quantization procedure for the activations is aligned with that utilized in BitNet. However, instead of adjusting activations prior to the application of non-linearities so that they lie within $[0, Q_b]$, we normalize all activations on a per-token basis to fall within $[-Q_b, Q_b]$, effectively eliminating the need for zero-point quantization. This adjustment streamlines both the coding process and the efficiency of system-level operations, while empirical results indicate it leads to minimal impact on model accuracy.

\paragraph{LLaMA-alike Components.}

LLaMA~\cite{llama, llama2} has become a standard model architecture for the open-source large language model ecosystem. To better serve the open-source community, the design of \text{BitNet b1.58} mirrors several LLaMA architectural choices. In particular, this model utilizes RMSNorm~\cite{rmsnorm}, SwiGLU~\cite{swiglu}, rotary embeddings~\cite{rope}, and excludes bias terms. This strategy allows \text{BitNet b1.58} to be readily compatible with widely used open-source libraries such as Huggingface, vLLM~\cite{vllm}, and llama.cpp\footnote{\url{https://github.com/ggerganov/llama.cpp}}, thus lowering integration barriers.

\section{Results}

We conducted a comparative analysis between \text{BitNet b1.58} and our implementation of the FP16 \text{LLaMA LLM} across multiple model sizes. To maintain experimental consistency, both models underwent pre-training on the RedPajama dataset~\cite{redpajama} for a total of 100 billion tokens. The zero-shot evaluation encompassed several language benchmarks, such as ARC-Easy~\cite{arc}, ARC-Challenge~\cite{arc}, Hellaswag~\cite{hellaswag}, Winogrande~\cite{winoGrande}, PIQA~\cite{piqa}, OpenbookQA~\cite{openbookqa}, and BoolQ~\cite{boolq}. Additionally, we provided validation perplexity scores on the WikiText2~\cite{wikitext2} and C4~\cite{c4} datasets.

For runtime assessment, we recorded GPU memory consumption and inference latency for both \text{LLaMA LLM} and \text{BitNet b1.58}. Measurements were obtained using the FasterTransformer\footnote{\url{https://github.com/NVIDIA/FasterTransformer}} library, which is specifically optimized for large language model inference on GPUs. The 2-bit kernel from Ladder~\cite{ladder} was incorporated to support \text{BitNet b1.58}. Inference cost was reported in terms of average time per output token, reflecting the primary expense during inference.

A summary of the models' perplexity and computational costs is presented in Table~\ref{tab:ppl}. The findings indicate that, in terms of perplexity, \text{BitNet b1.58} achieves parity with the full precision \text{LLaMA LLM} at the 3B parameter scale, while offering a 2.71-fold improvement in speed and using 3.55 times less GPU memory. Notably, the 3.9B variant of \text{BitNet b1.58} delivers 2.4 times faster performance and reduces memory usage by 3.32 times, in addition to substantially outperforming the 3B \text{LLaMA LLM}.

Detailed zero-shot accuracy for downstream tasks is detailed in Table~\ref{tab:zero-shot}. For these evaluations, we adopted the methodology from \emph{lm-evaluation-harness}\footnote{\url{https://github.com/EleutherAI/lm-evaluation-harness}}. The outcomes reveal that the performance differential between \text{BitNet b1.58} and \text{LLaMA LLM} diminishes as model scale increases. More critically, \text{BitNet b1.58} matches the accuracy of the full-precision baseline starting with the 3B size. Mirroring the trend observed in perplexity, \text{BitNet b1.58} 3.9B achieves superior performance over \text{LLaMA LLM} 3B while maintaining lower resource usage in terms of memory and latency. These results collectively establish \text{BitNet b1.58} as a Pareto improvement relative to leading LLM models.

\paragraph{Memory and Latency}

We expanded our experiments by increasing the model size to 7B, 13B, and 70B parameters, and analyzed the corresponding computational costs. As depicted in Figure~\ref{fig:latency-memory-energy}, both latency and memory usage trends reveal that larger models achieve greater acceleration. Notably, \text{BitNet b1.58} with 70B parameters demonstrates a 4.1$\times$ improvement in inference speed relative to the \text{LLaMA LLM} baseline. This performance gain is primarily attributed to the growth in computational time of \emph{nn.Linear} as models scale up. Similarly, memory requirements also increase with model size, though the contribution of embeddings, which are stored in full precision, diminishes proportionally in larger models. Latency and memory metrics were obtained using a 2-bit kernel, implying potential for further cost reductions through future optimization.

\paragraph{Energy}

We additionally assessed the estimated energy usage associated with arithmetic operations in both \text{BitNet b1.58} and \text{LLaMA LLM}. Our analysis centers on matrix multiplication, which is the primary contributor to computational expense in large language models. Figure~\ref{fig:energy} provides a breakdown of the respective energy profiles. For \text{BitNet b1.58}, the bulk of operations are INT8 additions, whereas \text{LLaMA LLM} heavily utilizes both FP16 additions and multiplications. Drawing on the energy estimation model from~\cite{energycost, pokebnn}, \text{BitNet b1.58} achieves a 71.4$\times$ reduction in arithmetic operation energy consumption for matrix multiplication when running on 7nm hardware. We also measured the end-to-end energy cost for models processing 512 tokens. The results indicate that as model sizes increase, \text{BitNet b1.58} exhibits progressively greater energy efficiency compared to the FP16 \text{LLaMA LLM} baseline. This increased efficiency stems from the fact that the proportion of computation spent on \emph{nn.Linear} layers rises with model size, while the impact of other components becomes less significant in larger architectures.

\paragraph{Throughput}

We evaluated the throughput of \text{BitNet b1.58} and the 70B parameter \text{LLaMA LLM} on two 80GB A100 GPUs, employing pipeline parallelism~\cite{gpipe} to enable \text{LLaMA LLM} 70B to fit onto the devices. For both models, the batch size was scaled up to the point at which GPU memory was fully utilized, with a fixed sequence length of 512. As presented in Table~\ref{tab:throughput}, \text{BitNet b1.58} 70B accommodates batch sizes up to 11 times larger than those possible with \text{LLaMA LLM}, which leads to an 8.9-fold increase in throughput.

\textbf{\text{BitNet b1.58} introduces a novel scaling relationship between model efficacy and inference expenditure.} Based on the empirical data in Figures~\ref{fig:latency-memory-energy} and \ref{fig:energy}, equivalences between various model capacities at 1.58-bit precision and at 16-bit precision are as follows:

\begin{itemize}
    \item The 13B BitNet b1.58 model demonstrates lower latency, reduced memory demands, and decreased energy consumption compared to a 3B FP16 LLM.
    \item The 30B BitNet b1.58 surpasses the 7B FP16 LLM in terms of latency, memory efficiency, and power usage.
    \item The 70B BitNet b1.58 model is more optimal than the 13B FP16 LLM across all three metrics: latency, memory, and energy consumption.
\end{itemize}

\paragraph{Training with 2T Tokens}

For large language models, the number of training tokens significantly impacts performance. To assess \text{BitNet b1.58}'s scalability with respect to token count, we trained a \text{BitNet b1.58} model on 2T tokens, adhering to the data preparation strategy from StableLM-3B~\cite{StableLM-3B-4E1T}, recognized as a leading open-source 3B model. Both \text{BitNet b1.58} and StableLM-3B were evaluated using a composite benchmark comprising Winogrande~\cite{winoGrande}, PIQA~\cite{piqa}, SciQ~\cite{sciq}, LAMBADA~\cite{lambada}, and ARC-easy~\cite{arc}. We present zero-shot accuracies in Table~\ref{tab:2t}, reporting the mean when both accuracy and normalized accuracy metrics are involved. Scores for StableLM 3b at 2T tokens are referenced directly from its technical documentation. Our experiments demonstrate that \text{BitNet b1.58} attains superior results across all evaluated tasks, implying that 1.58-bit LLMs maintain robust generalization capabilities.

\section{Discussion and Future Work}

\textbf{1-bit Mixture-of-Experts (MoE) LLMs}

The Mixture-of-Experts (MoE) paradigm continues to serve as an efficient solution for scaling Large Language Models (LLMs) cost-effectively. Although MoE architectures reduce floating-point operations substantially, they are often constrained by substantial memory demands and the bottleneck of inter-chip communication, both of which hinder practical deployment. Utilizing 1.58-bit LLMs presents a promising route to alleviate these constraints. Primarily, minimizing the model's memory requirements lessens the number of hardware units necessary for deployment. In addition, lower memory usage leads to less interconnect traffic due to smaller activation data, thereby minimizing communication overhead. Ideally, if all model parameters fit within a single processing chip, communication overhead could be eliminated entirely.

\textbf{Native Support for Long Sequences in LLMs}

With the advancement of LLM technology, efficiently handling longer sequences has become increasingly essential. The primary barrier to long-sequence inference is the memory required to store Key-Value (KV) caches. The introduction of \text{BitNet b1.58} marks progress toward inherent support for longer contexts by compressing activation precision from 16 bits down to 8 bits, which effectively permits doubling the manageable sequence length within identical hardware constraints. Future efforts may focus on further compressing these activations to as little as 4 bits, or potentially even less, through the adoption of 1.58-bit LLMs, thus expanding capacity without loss, which we propose as an avenue for continued research.

\textbf{LLMs on Edge and Mobile}

Deploying 1.58-bit LLMs stands to significantly enhance the applicability of large language models in environments with restricted computational resources, such as edge and mobile devices. Typically, these platforms face challenges related to limited memory and processing capabilities, restricting the complexity and scale of models that can be run. The notably lower memory and energy requirements of 1.58-bit LLMs make their deployment on such constrained devices practical, unlocking diverse use cases previously inaccessible on these platforms. In addition, since CPUs are the predominant processors found in edge and mobile systems, and since 1.58-bit LLMs exhibit favorable characteristics for CPU-based computation, \text{BitNet b1.58} can operate efficiently, boosting both performance and range of feasible applications on these devices.

\textbf{New Hardware for 1-bit LLMs}

Recent advancements—such as those demonstrated by Groq\footnote{\url{https://groq.com/}}—highlight the substantial promise of developing purpose-built hardware (e.g., Language Processing Units, or LPUs) tailored for LLM execution. Looking forward, we advocate for the creation and design of new hardware and computational infrastructures expressly optimized for 1-bit LLM architectures, building on the novel computational approach introduced in BitNet~\cite{bitnet}.